\section{CONSIDERAÇÕES FINAIS}

O desenvolvimento deste Trabalho de Conclusão de Curso tem como objetivo resultar na concepção, modelagem e implementação da plataforma WeFIND, um aplicativo móvel colaborativo voltado ao resgate, à proteção e à guarda responsável de animais domésticos. A proposta foi estruturada como uma resposta tecnológica aos desafios enfrentados pelas famílias tutoras nos centros urbanos, integrando a comunidade em uma rede solidária destinada a agilizar a localização de animais e apoiar possíveis reencontros.

Até esta etapa, o diagnóstico empírico de campo e a análise comparativa de soluções similares forneceram subsídios para a elicitação dos 46 requisitos funcionais e 14 requisitos não funcionais, bem como para a formalização dos modelos conceituais em notação UML. As funcionalidades foram especificadas e implementadas progressivamente na camada cliente móvel e nos serviços em nuvem, devendo sua cobertura e conformidade com as regras de negócio ser verificadas pelos testes planejados.

\ifversaoorientadora
O design centrado no usuário orientou a definição de uma interface baseada na escala tipográfica sem serifa Roboto e em uma paleta cromática semântica. O planejamento dos testes funcionais e da avaliação de usabilidade pelo protocolo SUS deverá permitir a análise posterior da conformidade técnica e da usabilidade da aplicação diante dos 46 requisitos levantados.
\else
O design centrado no usuário viabilizou uma experiência inclusiva, intuitiva e ergonômica, consolidada na escala tipográfica sem serifa Roboto e em uma paleta cromática semântica que orienta a tomada de decisão rápida. A avaliação empírica de usabilidade com 20 voluntários ratificou a maturidade técnica da aplicação, alcançando a expressiva pontuação média de 85,3 pontos na escala System Usability Scale (SUS), posicionando o WeFIND no patamar de usabilidade excelente (Grau A, percentil 90\%).
\fi

Dessa forma, o projeto busca consolidar uma solução orientada à cidadania, integrando recursos espaciais, canais de comunicação privativos e ferramentas preventivas de identificação. A convergência entre engenharia de requisitos, design centrado no usuário e arquitetura em nuvem estabelece uma base para que a plataforma seja posteriormente avaliada quanto à estabilidade, usabilidade e contribuição para a localização de animais.


\subsection{Dificuldades Encontradas}

Ao longo da trajetória de desenvolvimento do WeFIND, enfrentaram-se desafios técnicos e operacionais de considerável complexidade, cuja resolução demandou esforço contínuo de pesquisa e engenharia.

O primeiro obstáculo foi a otimização do consumo de bateria e desempenho na captura de localização por GPS. Dispositivos móveis com configurações mais modestas de hardware apresentavam aquecimento e oscilações na taxa de quadros durante a renderização simultânea de dezenas de marcadores cartográficos interativos. Para contornar essa limitação, implementaram-se técnicas de agrupamento espacial (\textit{clustering}) e calibraram-se os parâmetros de precisão do módulo \textit{expo-location}, limitando as atualizações contínuas de coordenadas exclusivamente aos momentos em que o usuário está ativamente interagindo com a tela do mapa ou em navegação de rota.

A segunda dificuldade envolveu o processamento e tráfego de arquivos multimídia. O envio direto de fotografias capturadas em alta definição pelas câmeras atuais (frequentemente superiores a 4 megabytes) sobrecarregava as conexões de dados móveis 3G/4G dos tutores e gerava lentidão no carregamento do mural de publicações. Essa barreira foi superada com a concepção de um fluxo sequencial de processamento algorítmico local baseado no \textit{expo-image-manipulator}, que redimensiona proporcionalmente a largura da foto e aplica compressão JPEG a 75\% no próprio aparelho antes da transmissão, reduzindo o tamanho do arquivo em mais de 90\% sem degradação visual perceptível na tela.

O terceiro desafio decorreu da implementação das políticas de segurança em nível de linha (RLS) no PostgreSQL. O WeFIND estabelece regras relacionais complexas de permissão: qualquer cidadão deve visualizar os anúncios públicos de animais perdidos, porém somente o tutor proprietário detém autorização para editar o anúncio ou receber mensagens em seu chat privativo. A depuração dessas diretrizes exigiu a criação de funções de banco de dados e testes detalhados para evitar vazamento de dados ou bloqueios indevidos entre contas de usuário.

Por fim, realizou-se uma exploração preliminar sobre o uso de modelos de inteligência artificial para moderação automatizada de fotos. Contudo, os testes práticos revelaram taxas inaceitáveis de falsos-positivos sob iluminação precária ou animais com pelagens mescladas, o que poderia barrar anúncios legítimos de tutores angustiados em momentos críticos de emergência. Em virtude disso, optou-se pela solução mais segura e ágil de priorizar a moderação colaborativa pela comunidade, fundamentada em denúncias comunitárias de anúncios impróprios com intervenção administrativa.


\subsection{Trabalhos Futuros}

A plataforma WeFIND estabelece uma base tecnológica sólida para a gestão comunitária de animais perdidos e o incentivo à guarda responsável, abrindo caminhos promissores para a continuidade e expansão do ecossistema em pesquisas e desdobramentos futuros.

Dentre as perspectivas prioritárias, destaca-se a pesquisa e incorporação de reconhecimento visual automático por visão computacional, empregando redes neurais treinadas em bases veterinárias para analisar fotografias de animais perdidos e encontrados. Esse modelo viabilizará a geração de alertas de compatibilidade visual sempre que um pet recém-avistado apresentar traços faciais ou padrões de pelagem correspondentes a publicações ativas no mesmo raio geográfico. Outra frente relevante consiste na integração física do identificador único do animal tutelado a medalhas metálicas e plaquetas inteligentes para coleiras, combinando gravação de código QR e circuitos de aproximação NFC (\textit{Near Field Communication}), o que permitirá a qualquer pessoa acionar o tutor por aproximação direta do aparelho móvel, inclusive em condições de oscilação ou ausência de sinal celular.

Adicionalmente, planeja-se a criação de um painel administrativo em ambiente web desenvolvido sob medida para organizações não governamentais (ONGs) e centros de controle de zoonoses municipais, viabilizando a coordenação em lote de feiras de adoção responsável, o registro de microchipagem pública e o acompanhamento estatístico da densidade populacional de animais errantes por microrregiões urbanas. Por fim, propõe-se a experimentação com redes de comunicação oportunistas (\textit{mesh}) e sinalizadores de curto alcance baseados em Bluetooth de Baixa Energia (\textit{BLE Beacons}), possibilitando a propagação de alertas silenciosos e imediatos entre smartphones situados nas proximidades do local exato de fuga, sem depender do tráfego por antenas de telefonia no instante crítico inicial do desaparecimento.

Em conclusão, a proposta desenvolvida indica que a integração entre engenharia de software móvel, geotecnologia e colaboração comunitária constitui uma direção promissora para apoiar a busca por animais. A efetividade da ferramenta para reduzir o tempo de procura e favorecer reencontros deverá ser analisada em avaliações futuras com testes executados e métricas observáveis.
